\section{Related Work}
\label{sec:related}

% B1. Signal models, modal analysis, and differentiable percussion models.
\paragraph{Analysis and synthesis of percussive sound.}
Signal models of drum sounds predate deep learning, as in the drum variant of the Karplus--Strong algorithm, which turns a noise burst into an unpitched decaying sound \citep{karplus1983digital}.
Modal approaches instead describe a drum as a set of decaying modes derived from recordings or measurements.
\citet{kirby2021evolution} modeled the pitch glide of each mode of a tom-tom across strike intensities, and \citet{nagata2025sound} synthesized drum sounds from transfer functions obtained by experimental modal analysis.
These models reproduce a particular instrument closely, but they are built from that instrument rather than learned across a heterogeneous sample library.
Differentiable modal synthesis learns such parameters from data, for example by predicting differentiable modal resonators from the shape, material, and strike position of an object \citep{diaz2023rigid}.
Learning the frequencies themselves is harder, because the loss surface over sinusoidal frequency is non-convex and contains many local minima \citep{hayes2023sinusoidal}.
In DDSP-Piano, a second training phase that fine-tunes the inharmonicity and detuning parameters, among others, did not improve the perceived quality in a listening test \citep{renault2023ddsp}.
In DrumBlender, the sinusoidal frequencies are therefore not learned but extracted before training by peak tracking on a constant-Q transform, and the noise and transient parameters are predicted by encoders adapted from the convolutional encoder of the SoundStream codec \citep{shier2023differentiable,zeghidour2022soundstream}.
For stochastic components, NoiseBandNet predicts the time-varying amplitudes of a bank of narrow noise bands \citep{barahonarios2024noisebandnet}.
In contrast, this work applies a hybrid constant-Q and short-time Fourier analysis to every instrument class in a sample library, keeps the analyzed tracks and their initial phases, and learns only bounded corrections of their frequencies and amplitudes.
Its noise bands follow NoiseBandNet, but their amplitudes are predicted jointly with the modal corrections and conditioned on them, and its transient component restricts the modal-driven TCN of DrumBlender to the modal waveform and a single clip-level code.

% B2. Generative drum models and the controls they expose.
\paragraph{Neural drum synthesis and its controls.}
Beyond the descriptor-conditioned models discussed in \sectionref{sec:intro}, neural drum synthesis has explored several generative families and forms of control.
\citet{drysdale2020adversarial} conditioned a WaveGAN-based model on the drum class and explored its latent space by interpolation.
StyleWaveGAN adapts a style-based generator to raw waveforms and combines discrete drum-type control with continuous perceptual descriptors and per-class amplitude envelopes \citep{lavault2022stylewavegan}.
CRASH applies score-based diffusion directly to raw audio and can regenerate an existing sound from its latent code under a deterministic sampler \citep{rouard2021crash}.
In a user study of a style-based drum generator, unsupervised latent directions were rated slightly higher than random and supervised ones, although not significantly \citep{ramires2022study}.
\citet{drysdale2024latent} explored the latent space of a drum generator through latent arithmetic, dimensionality reduction, and visualization.
\citet{waissbluth2022synthesis} instead reconstructed snare drum samples as sparse mixtures of a small fixed sample library with a variational autoencoder, so that moving through its latent space swaps samples in and out of the mixture.
In these systems, an existing sample can be edited only through the generator's own representation of it, obtained by an encoder or by inversion \citep{nistal2022drumganvst,rouard2021crash}, so the edit applies to the generator's approximation of the sample.
In contrast, this work starts from the given recording, reconstructs it with modal, noise, and transient components whose parameters are derived from an analysis of that recording rather than from a fixed library, and applies the controls to that reconstruction.

% B3. Learned latent spaces and controls without labels.
\paragraph{Label-free control of learned audio representations.}
Neural audio codecs learn compact latent representations for high-fidelity reconstruction \citep{zeghidour2022soundstream,kumar2023high}, but their latent dimensions are not trained to correspond to perceptual attributes.
Variational autoencoders impose a prior that makes the latent space easier to traverse, and \citet{esling2018generative} further aligned latent distances between instrument classes with perceptual timbre spaces, which requires instrument labels and human dissimilarity ratings.
Sparse autoencoders offer a post-hoc alternative that reconstructs the activations of a trained model through a wide layer of sparsely active features \citep{gao2025scaling}.
Such features have been identified in audio latent spaces \citep{paek2025learning} and used to steer large generative music models \citep{singh2026discovering}.
For percussion, \citet{shier2024realtime} learned without parameter supervision to map onset features of acoustic snare hits to parameter changes of a differentiable TR-808-style snare synthesizer.
\citet{shier2025designing} extended this timbre remapping to non-differentiable synthesizers and reported that the temporal morphology of a hit, namely how its loudness and spectrum evolve over time, was poorly mapped.
In contrast, this work fits a sparse autoencoder post hoc to the shared representation of a frozen model that reconstructs a given sample, so that label-free edits reach the modal and noise components through their synthesis parameters.
